\section{Hardware and Operating System Discussion}
\label{sec:hardware-os-discussion}

The tested and referenced hardware used in this project includes the following:

\begin{itemize}
\item \textbf{Apple M3 Max, 48~GB unified memory, macOS, ARM64:} used for the
matrix v1--v7 campaign, the v6 four-workload campaign, and one v8 matrix group.
\item \textbf{MacBook Pro with Apple M2 Pro, 32~GB memory, macOS, ARM64:}
v8.0.0 development, integration and unit testing, custom metric analysis
\item \textbf{AMD Ryzen~9~7900X, 32~GB RAM, Ubuntu~24.04 under WSL, x86-64:}
used for the v7 four-workload suite.
\item \textbf{Intel i7-1165G7, 16~GB RAM, Ubuntu~24.04.3, x86-64:} used for
Jacob's v2 measurements and the newer v8 matrix runs.
\item \textbf{Intel i7-11370H, 16~GB RAM, WSL and Linux Mint, x86-64:} used for
the pthread comparison runs.
\item \textbf{Additional team systems:} newer Apple Silicon machines and other
Intel and Ryzen hosts appeared in the shared team inventory, although not all
of them contributed balanced measurements across allocator versions and
workloads.
\end{itemize}

These hardware differences matter because allocator behavior is tied closely to cache hierarchy, memory bandwidth,
scheduler behavior, and runtime overheads from the operating system environment. The Apple M3 Max results
represent an ARM64 macOS platform with high memory bandwidth and unified memory, while the Ryzen and
Intel hosts represent x86-64 Linux or WSL environments with different cache organizations, thread scheduling
behavior, and background-system costs. For that reason, environment should be treated as a source of
experimental variation rather than noise to be ignored. It cannot be assumed that a result that holds on one
machine will behave identically on every processor and operating system and should instead be interpreted with
the knowledge that the results may have been influenced by the machine they were run on.
